\chapter{Evaluation}
\label{ch:evaluation}

\section{Effektivität: Retrieval-Qualität}

Zur Bewertung der Retrieval-Qualität wurden zwei manuell
zusammengestellte Evaluationsdatensätze verwendet. Der reguläre
Evaluationsdatensatz umfasst 20 natürlichsprachliche Suchanfragen aus
verschiedenen Themenbereichen, darunter Geschichte, Politik und Gesellschaft,
Biologie und Medizin, Naturwissenschaften, Geografie, Technologie und Kultur.
Für jede Anfrage wurde ein erwarteter Wikipedia-Artikel als Ground Truth
hinterlegt.

Ergänzend wurde ein separater Edge-Case-Datensatz erstellt. Dieser enthält
gezielt Anfragen, bei denen besondere Schwierigkeiten für das Retrieval zu
erwarten sind. Dazu gehören ähnlich benannte Entitäten, Aliase beziehungsweise
alternative Bezeichnungen, Tippfehler, sehr kurze Suchanfragen sowie eine
Anfrage zu einem Thema außerhalb des verwendeten Korpus. Die Edge Cases wurden getrennt
vom regulären Evaluationsdatensatz betrachtet, da sie nicht die allgemeine
Retrieval-Qualität abbilden, sondern gezielt das Verhalten des Systems in
problematischen Situationen untersuchen.

Als Retrieval-Ergebnis werden die ersten $k$ unterschiedlichen
Wikipedia-Artikel betrachtet. Da ein Artikel im Vektorindex durch mehrere
Text-Chunks repräsentiert werden kann, werden mehrere Treffer desselben
Artikels bei der Evaluation zu einem Artikeltreffer zusammengefasst.

Zur quantitativen Bewertung werden Recall@k, Precision@k, nDCG@k und der
Mean Reciprocal Rank (MRR) verwendet. Recall@k beschreibt den Anteil der
relevanten Artikel, die innerhalb der ersten $k$ Ergebnisse gefunden werden.
Precision@k misst den Anteil relevanter Treffer unter den ersten $k$
Ergebnissen. nDCG@k berücksichtigt zusätzlich die Position eines relevanten
Treffers und bewertet früher platzierte Treffer höher. Der Reciprocal Rank
entspricht dem Kehrwert des Rangs des ersten relevanten Treffers; der MRR
bildet den Mittelwert dieses Wertes über alle Evaluationsanfragen.

Formal sei $G$ die Menge der relevanten Ground-Truth-Artikel und
$R_k$ die Menge der ersten $k$ Retrieval-Ergebnisse. Dann ist Recall@k
definiert als

\[
\mathrm{Recall@k} = \frac{|R_k \cap G|}{|G|},
\]

während Precision@k durch

\[
\mathrm{Precision@k} = \frac{|R_k \cap G|}{k}
\]

\par\smallskip
\noindent\begin{minipage}[t]{\linewidth}
gegeben ist. Für nDCG@k wird die Position relevanter Treffer
berücksichtigt. Bei binärer Relevanz wird zunächst

\[
\mathrm{DCG@k} =
\sum_{i=1}^{k}
\frac{\mathrm{rel}_i}{\log_2(i+1)}
\]
\end{minipage}\par\smallskip

berechnet und anschließend durch den idealen DCG-Wert normalisiert:

\[
\mathrm{nDCG@k} =
\frac{\mathrm{DCG@k}}{\mathrm{IDCG@k}}.
\]

Der Reciprocal Rank einer Anfrage ist

\[
\mathrm{RR} = \frac{1}{r},
\]

wobei $r$ dem Rang des ersten relevanten Treffers entspricht. Wird kein
relevanter Artikel gefunden, ist der Wert 0. Der Mean Reciprocal Rank ergibt
sich als Mittelwert der Reciprocal-Rank-Werte über alle Anfragen.

Die Metriken betrachten unterschiedliche Aspekte der Retrieval-Qualität.
Recall@k zeigt, ob die relevanten Artikel innerhalb der betrachteten
Ergebnisse gefunden werden, während Precision@k den Anteil relevanter
Treffer innerhalb dieser Ergebnisse beschreibt. nDCG@k und MRR
berücksichtigen zusätzlich die Rangposition und unterscheiden damit
zwischen relevanten Treffern, die früh oder erst später in der
Ergebnisliste erscheinen.

\section{Effektivität: Ergebnisse}

Für den regulären Evaluationsdatensatz wurde das Retrieval zunächst mit
$k=5$ und anschließend mit $k=10$ ausgewertet. Bei beiden Einstellungen
wurden für alle 20 Anfragen die jeweils als relevant definierten
Wikipedia-Artikel innerhalb der betrachteten Treffer gefunden. Damit wurde
sowohl für $k=5$ als auch für $k=10$ ein mittlerer Recall von $1{,}000$
erreicht. Der MRR betrug in beiden Fällen $0{,}950$, der mittlere nDCG-Wert
$0{,}963$.

\begin{table}[ht]
    \centering
    \caption{Retrieval-Ergebnisse des regulären Evaluationsdatensatzes}
    \label{tab:retrieval-regulaer}
    \begin{tabular}{lrrrr}
        \toprule
        Einstellung & Recall@k & Precision@k & nDCG@k & MRR \\
        \midrule
        $k=5$  & 1,000 & 0,200 & 0,963 & 0,950 \\
        $k=10$ & 1,000 & 0,100 & 0,963 & 0,950 \\
        \bottomrule
    \end{tabular}
\end{table}

Die mittlere Precision betrug bei $k=5$ $0{,}200$ und bei $k=10$
$0{,}100$. Dieser Rückgang ist durch die Struktur des Evaluationsdatensatzes
zu erklären: Für jede der 20 Anfragen ist genau ein relevanter
Ground-Truth-Artikel definiert. Wird dieser gefunden, beträgt die
Precision daher bei fünf betrachteten Ergebnissen $1/5$ und bei zehn
betrachteten Ergebnissen $1/10$. Die Erhöhung von $k$ von 5 auf 10 führte
somit auf diesem Datensatz zu keiner Verbesserung von Recall, nDCG oder MRR.

Zusätzlich wurde das Retrieval auf dem separaten Edge-Case-Datensatz mit
zwölf Anfragen bei $k=5$ untersucht. Ohne Query Analyzer wurden ein
mittlerer Recall@5 von $0{,}833$, eine Precision@5 von $0{,}167$, ein
nDCG@5 von $0{,}803$ und ein MRR von $0{,}792$ erreicht. Mit aktiviertem Query Analyzer stiegen diese Werte auf $0{,}917$, $0{,}183$,
$0{,}886$ beziehungsweise $0{,}875$. Gegenüber der Baseline entspricht dies
einer absoluten Verbesserung von $0{,}083$ bei Recall@5, nDCG@5 und MRR sowie
von $0{,}017$ bei Precision@5.

\begin{table}[ht]
    \centering
    \caption{Retrieval-Ergebnisse auf dem Edge-Case-Datensatz bei $k=5$}
    \label{tab:retrieval-edge-cases}
    \begin{tabular}{lrrrr}
        \toprule
        Einstellung & Recall@5 & Precision@5 & nDCG@5 & MRR \\
        \midrule
        Ohne Query Analyzer & 0,833 & 0,167 & 0,803 & 0,792 \\
        Mit Query Analyzer  & 0,917 & 0,183 & 0,886 & 0,875 \\
        \bottomrule
    \end{tabular}
\end{table}

Die Verbesserung ist in diesem Versuch auf die Anfrage
\texttt{Cathlic Chuch} zurückzuführen. Ohne Vorverarbeitung wurde der
erwartete Artikel innerhalb der ersten fünf Ergebnisse nicht gefunden.
Der Query Analyzer normalisierte diese Anfrage zu
\texttt{Catholic Church}; anschließend wurde der erwartete Artikel auf
Rang 1 gefunden. Andere untersuchte Edge Cases, beispielsweise
\texttt{Who was George IV?}, wurden durch diese Normalisierung nicht
verändert.

Der Edge-Case-Datensatz enthält außerdem eine als Out-of-Scope markierte
Anfrage zu Albert Einstein, für die kein relevanter Ground-Truth-Artikel
definiert ist. Die aktuelle Implementierung der klassischen
Retrieval-Metriken weist einer Anfrage ohne relevante Ground-Truth-Titel
den Wert 0 zu. Dieser Fall ist daher nicht als gewöhnlicher
Retrieval-Fehler zu interpretieren und muss bei der Einordnung der
aggregierten Edge-Case-Werte gesondert berücksichtigt werden.

\section{Effizienz: Indexierung}

Die Ressourcennutzung während eines Indexierungslaufs auf dem Raspberry Pi
wurde in Intervallen von zehn Sekunden aufgezeichnet. Erfasst wurden unter
anderem die systemweite Speichernutzung, der RSS-Speicherbedarf des
WikiPod-Prozesses und von OpenSearch, die Swap-Nutzung sowie CPU-Temperatur
und Systemlast.

Während der aufgezeichneten Messung lag die maximale systemweite
Speichernutzung bei $6224{,}3$ MB. Für den WikiPod-Prozess wurde ein
maximaler RSS-Wert von $2871{,}1$ MB und für OpenSearch ein maximaler
RSS-Wert von $4731{,}6$ MB gemessen. Die beiden RSS-Maximalwerte traten
nicht notwendigerweise zum selben Zeitpunkt auf und werden daher nicht
addiert. Während der gesamten Messung wurde kein Swap-Speicher verwendet.
Die höchste gemessene CPU-Temperatur betrug $71{,}4\,^\circ$C.

Die Monitoring-Aufzeichnung umfasst insgesamt 62 Messpunkte über einen
Zeitraum von etwa $10{,}2$ Minuten. Dieser Zeitraum ist nicht mit der
exakten Indexierungsdauer gleichzusetzen, da das Monitoring vor dem
eigentlichen Lauf gestartet und nach dessen Ende beendet wurde.

Ergänzend wurde die Indexierungsdauer auf dem Raspberry Pi in einem
separaten Lauf direkt als Wall-Clock-Zeit gemessen. Hierfür wurde derselbe
Top-100-ZIM-Korpus mit 194 Artikeln verwendet. Bei einem konfigurierten
Speicherbudget von $500$ MB wurden $61{,}6$ MB Inhalt ausgewählt und
insgesamt 3875 Text-Chunks erzeugt und in einen separaten OpenSearch-Index
geschrieben. Die gesamte Indexierung dauerte $606{,}76$ Sekunden, entsprechend
etwa $10$ Minuten und $7$ Sekunden.

Ergänzend wurde die Indexierung des für die Retrieval-Evaluation verwendeten
Top-100-ZIM-Korpus lokal reproduziert. Aus dem Korpus wurden 194 von 194
verfügbaren Artikeln ausgewählt. Der dabei berücksichtigte Inhalt umfasste
$61{,}6$ MB bei einem konfigurierten Speicherbudget von $500$ MB. Aus den
Artikeln wurden insgesamt 3875 Text-Chunks erzeugt und in OpenSearch
indexiert. Der resultierende OpenSearch-Index enthielt entsprechend 3875
Dokumente und belegte $42{,}1$ MB Speicherplatz.

Dieser lokale Lauf dient der Reproduzierbarkeit der für die
Retrieval-Evaluation verwendeten Indexbasis. Die zuvor beschriebenen
Ressourcenmessungen stammen dagegen vom Raspberry Pi und werden daher nicht
direkt mit den lokal gemessenen Werten vermischt.


\section{Effizienz: Anfragebearbeitung}

Die Latenz der Retrieval-Stufe wurde zusätzlich in einem lokalen
Reproduktionslauf mit den 20 Anfragen des regulären Evaluationsdatensatzes
und $k=5$ gemessen. Dabei wurde jeweils der vollständige Aufruf der
Retrieval-Komponente erfasst, der sowohl die Berechnung des Query-Embeddings
als auch die anschließende $k$-NN-Suche in OpenSearch umfasst. Ein
vorangestellter Warm-up-Aufruf wurde nicht in die Messwerte einbezogen.

Die mittlere Retrieval-Latenz betrug $27{,}13$ ms pro Anfrage, bei einem
Median von $28{,}73$ ms. Die gemessenen Einzelwerte lagen zwischen
$15{,}99$ ms und $41{,}26$ ms. Da diese Messung auf dem lokalen
Entwicklungssystem und nicht auf dem Raspberry Pi durchgeführt wurde,
dienen die Werte der Charakterisierung des Evaluationsaufbaus und sind
nicht als Laufzeitmessung der Zielhardware zu interpretieren.

Die gleiche Retrieval-Messung wurde anschließend auf dem Raspberry Pi
wiederholt. Verwendet wurden ebenfalls die 20 Anfragen des regulären
Evaluationsdatensatzes, $k=5$ und ein vorangestellter Warm-up-Aufruf, der
nicht in die Messwerte einbezogen wurde. Gemessen wurde wiederum der
vollständige Aufruf der Retrieval-Komponente einschließlich
Query-Embedding und $k$-NN-Suche in OpenSearch.

Auf dem Raspberry Pi betrug die mittlere Retrieval-Latenz $62{,}78$ ms
pro Anfrage bei einem Median von $53{,}76$ ms. Die Einzelwerte lagen
zwischen $46{,}40$ ms und $115{,}56$ ms. Damit lag die mittlere
Retrieval-Latenz in diesem Versuchsaufbau auf dem Raspberry Pi etwa beim
$2{,}31$-Fachen des lokal gemessenen Mittelwerts von $27{,}13$ ms.

Zur Untersuchung der Generierungslatenz wurden auf dem Raspberry Pi drei
lokale Sprachmodelle im GGUF-Format verglichen. Für jedes Modell wurden
dieselben 20 Anfragen des regulären Evaluationsdatensatzes verwendet.
Das Retrieval wurde dabei einmal pro Anfrage ausgeführt und für alle
Modelle wiederverwendet. Die gemessene Zeit umfasst daher ausschließlich
den Aufruf der Generierung und nicht die Retrieval-Latenz.

\begin{table}[ht]
    \centering
    \caption{Mittlere Generierungslatenz lokaler Sprachmodelle auf dem Raspberry Pi}
    \label{tab:pi-generation-latency}
    \begin{tabular}{lrr}
        \toprule
        Modell & Modellgröße & Mittlere Latenz \\
        \midrule
        Llama 3.2 1B Q4 & 771 MB & 29,18 s \\
        Phi-3.5 Mini Q4 & 2,3 GB & 152,33 s \\
        Qwen2.5 1.5B Q4 & 1,1 GB & 27,97 s \\
        \bottomrule
    \end{tabular}
\end{table}

In diesem Versuchsaufbau erreichte Qwen2.5 1.5B mit $27{,}97$ Sekunden
die niedrigste mittlere Generierungslatenz, knapp vor Llama 3.2 1B mit
$29{,}18$ Sekunden. Phi-3.5 Mini benötigte mit $152{,}33$ Sekunden
deutlich mehr Zeit. Die Messung bewertet ausschließlich die Laufzeit der
Generierung; aus diesen Werten kann keine Aussage über die Qualität der
generierten Antworten abgeleitet werden.

Ergänzend wurde die Ressourcennutzung während einer einzelnen
Generierung mit Qwen2.5 1.5B auf dem Raspberry Pi in Intervallen von zwei
Sekunden aufgezeichnet. Die Generierung dauerte in diesem Lauf
$27{,}42$ Sekunden. Die systemweite Speichernutzung stieg von
$4577{,}8$ MB vor der Generierung auf maximal $5770{,}7$ MB und lag am
Ende der Messung bei $4593{,}5$ MB. Gegenüber dem Ausgangswert entspricht
dies einem Anstieg bis zum Maximum von etwa $1192{,}9$ MB.

Die Swap-Belegung lag bereits vor Beginn der Generierung bei etwa
$2048$ MB. Während des Messzeitraums wurden zwei zusätzliche
Swap-In-Seiten und keine Swap-Out-Seiten registriert. Die bereits
vorhandene Swap-Belegung kann daher nicht der untersuchten Generierung
zugeschrieben werden. Die CPU-Temperatur stieg während der Messung von
$44{,}4\,^\circ$C auf maximal $56{,}0\,^\circ$C und lag am Ende bei
$46{,}1\,^\circ$C.

Der separat erfasste RSS-Wert des WikiPod-Prozesses bildet in diesem Lauf
den Speicherbedarf der eigentlichen Modellgenerierung nicht verlässlich
ab. Aus diesem Grund wird für diese Messung nur die Veränderung der
systemweiten Speichernutzung berichtet und nicht daraus auf den
prozessspezifischen Speicherbedarf des Sprachmodells geschlossen.

\section{Einordnung der Ergebnisse}

Die Ergebnisse zeigen für den regulären Evaluationsdatensatz eine hohe
Retrieval-Qualität. Die Aussagekraft ist jedoch durch den Umfang des
Evaluationsaufbaus begrenzt. Der reguläre Datensatz umfasst lediglich
20 Anfragen und definiert für jede Anfrage genau einen relevanten
Ground-Truth-Artikel. Die gemessenen Werte sind daher als Ergebnis dieses
konkreten Evaluationsdatensatzes zu verstehen und erlauben keine allgemeine
Aussage über beliebige Wikipedia-Anfragen.

Auch der verwendete Korpus ist begrenzt. Für die reproduzierte
Retrieval-Evaluation wurde eine Top-100-ZIM-Datei mit einer Dateigröße von
318 MB verwendet. Beim Einlesen wurden daraus 194 Artikel verarbeitet und
3875 Chunks erzeugt. Die Ergebnisse können deshalb nicht unmittelbar auf
einen wesentlich größeren Wikipedia-Korpus übertragen werden.

Der Edge-Case-Datensatz verdeutlicht zudem, dass die aggregierten Metriken
von der Definition der Ground Truth abhängen. Insbesondere erhält die
Out-of-Scope-Anfrage ohne relevante Ground-Truth-Titel in der aktuellen
Implementierung für die klassischen Retrieval-Metriken jeweils den Wert
0. Sie muss deshalb getrennt von gewöhnlichen Retrieval-Fehlern
interpretiert werden.

\begingroup
\emergencystretch=1em
Der zusätzlich eingesetzte LLM-Judge erwies sich in der aktuellen
Konfiguration nur eingeschränkt als quantitative Bewertungsmetrik. In einem
frischen Lauf auf dem Edge-Case-Datensatz wurden insgesamt 22
Retrieval-Ergebnisse bewertet. Dabei vergab der Judge viermal den Wert 0, zweimal den
Wert 1, fünfzehnmal den Wert 2 und einmal den Wert 3 auf einer Skala von 0
bis 3.\par
\endgroup

Die Einzelbewertungen zeigen jedoch mehrere Inkonsistenzen. So wurde bei der
Anfrage \texttt{Who was George IV?} auch ein Treffer zum Artikel
\texttt{George III} mit 2 von 3 Punkten bewertet. Bei der Anfrage \texttt{Cathlic Chuch} wurde der nach der
Retrieval-Normalisierung korrekt gefundene Artikel \texttt{Catholic Church}
dagegen mit 0 Punkten bewertet. Der Judge erhielt dabei die ursprüngliche
Anfrage und nicht die für das Retrieval normalisierte Form. Auch bei der als
Out-of-Scope definierten Anfrage zu Albert Einstein erhielt ein Treffer zum
Artikel \texttt{Martin Luther King Jr.} 2 Punkte. Die Judge-Bewertungen werden
daher nicht als gleichwertige quantitative Qualitätsmetrik neben Recall,
nDCG und MRR verwendet, sondern als experimentelle Ergänzung betrachtet.

Die Messungen auf dem Raspberry Pi zeigen zudem deutliche Unterschiede
zwischen Retrieval und Antwortgenerierung hinsichtlich der Laufzeit. Die
mittlere Retrieval-Latenz lag mit $62{,}78$ ms deutlich unter einer Sekunde,
während die untersuchten lokalen Sprachmodelle für die Generierung einer
Antwort im Mittel zwischen $27{,}97$ s und $152{,}33$ s benötigten. In dem
untersuchten Aufbau stellt damit die lokale Sprachmodellgenerierung den
zeitlich dominierenden Teil der Anfragebearbeitung dar.

Auch zwischen den untersuchten Sprachmodellen bestanden deutliche
Laufzeitunterschiede. Qwen2.5 1.5B erreichte mit $27{,}97$ s die niedrigste
mittlere Generierungslatenz, gefolgt von Llama 3.2 1B mit $29{,}18$ s.
Phi-3.5 Mini benötigte dagegen im Mittel $152{,}33$ s. Da in diesem
Modellvergleich keine quantitative Qualitätsmetrik für die generierten
Antworten erhoben wurde, lässt sich aus der geringeren Latenz allein keine
Aussage über die Eignung eines Modells hinsichtlich der Antwortqualität
ableiten.

Die Ressourcenmessung einer einzelnen Qwen2.5-Generierung zeigte einen
Anstieg der systemweiten Speichernutzung um maximal etwa $1{,}19$ GB
gegenüber dem Ausgangswert. Dabei wurde keine relevante zusätzliche
Swap-Out-Aktivität beobachtet. Da der prozessspezifische Speicherbedarf
der Modellgenerierung in dieser Messung nicht verlässlich erfasst wurde,
ist dieser Wert als Veränderung der gesamten Systemspeichernutzung und
nicht als direkter Speicherbedarf des Modells zu interpretieren.